% \vspace{-3mm}
\section{Related Work}
% \vspace{-2mm}

\noindent\textbf{Video generation.}
Diffusion-based generative models~\cite{ho2020denoising, lipman2022flow, song2020score} have become a leading approach for video generation.
Built upon latent diffusion models~\cite{rombach2022high}, recent methods~\cite{chen2024videocrafter2, guo2023animatediff, yang2024cogvideox} perform generative modeling in the latent space, enabling more efficient synthesis of high-quality videos.
In parallel, autoregressive video generation approaches~\cite{chen2024diffusion, henschel2025streamingt2v, kim2024fifo} have attracted increasing attention for their potential to extend generation to arbitrary-length sequences, thereby offering a promising basis for building world models.
Driven by scalable architectures~\cite{peebles2023scalable} and increasingly sophisticated data curation pipelines, large-scale video generation systems~\cite{veo, sora, wan, kling, gao2025seedance, kong2024hunyuanvideo} trained on large-scale datasets have further exhibited emergent zero-shot abilities to understand, model, and manipulate the visual world~\cite{wiedemer2025video}, bringing physical world simulation closer to reality.

\noindent\textbf{Video world models.}
With the development of video diffusion transformers and the rapid advances in text-to-video and image-to-video generation, a growing body of world-modeling research has begun to repurpose video diffusion models as visual simulators~\cite{yang2023unisim,bar2025navigation,liang2024dreamitate}.
Rather than learning compact abstract states for downstream planning and decision making, these methods directly model future visual observations through generative video prediction.
We therefore refer to this family of approaches as video world models, highlighting their central reliance on video generation for simulating the evolution of the visual world.
Recent video world models have been increasingly studied in embodied scenarios, where agents must anticipate the visual consequences of actions and interactions.
Representative applications span robotics~\cite{mereu2025generative,yang2023learning,li2025unified,li2026causal,ye2026world,gao2026dreamdojo}, video games~\cite{alonso2024diffusion,he2025matrix,valevski2024diffusion,xiao2025worldmem, fan2022minedojo}, autonomous driving~\cite{agarwal2025cosmos,hu2023gaia}, and physical simulation~\cite{li2025pisa,yuan2026inference}, suggesting their potential as general-purpose simulators for dynamic visual environments.
Beyond open-loop video rollout, recent efforts~\cite{hyworld2025,worldplay2025,yu2025cam} further move toward interactive and temporally consistent world modeling, where the simulated environment can react to user or agent actions while preserving coherent scene structure, object dynamics, and long-range temporal consistency.

\noindent\textbf{Video diffusion model distillation.}
Real-time generation is a crucial requirement for building world models.
A common strategy for video diffusion models is to accelerate sampling through distillation~\cite{salimans2022progressive, geng2025mean, frans2024one}, which enables few-step inference and thus supports real-time synthesis.
Several works~\cite{sauer2024fast, sauer2024adversarial, kang2024distilling, lin2025diffusion, lin2024sdxl, lin2025autoregressive} introduce adversarial training objectives to reduce the number of denoising steps, but such methods are often prone to optimization instability and model collapse.
Another line of approaches~\cite{yin2024one, yin2024improved, lu2025adversarial} leverages variational score distillation~\cite{wang2023prolificdreamer} to obtain strong few-step generation performance.
Beyond these general acceleration techniques, CausVid~\cite{yin2025slow} tackles real-time autoregressive generation by transferring knowledge from a bidirectional diffusion teacher to a causal student model.
Building on this causal generation setup, Self-Forcing~\cite{huang2025self} further improves rollout training to reduce the exposure bias accumulated during long-horizon generation.
